\chapter{Domain integration}
\label{chap:domain-integration}

\begin{quote}
Spec dir:
\passthrough{\lstinline!specs/2026-05-02-phase-3-domain-integration/!}.
Version: \passthrough{\lstinline!0.3.0!}.
\end{quote}

\subsection{What this gives you}\label{what-this-gives-you}

A primitive that turns a scalar \passthrough{\lstinline!Field<double>!}
into a single number --- the discrete integral

\[
\int_{\Omega} f(\mathbf{x})\, dV \;\approx\; h_x h_y h_z \sum_{i,j,k} f_{ijk}
\]

over the periodic Cartesian domain
\passthrough{\lstinline![0, Nx*hx) × [0, Ny*hy) × [0, Nz*hz)!}.

This is the integration backbone the upcoming functional-evaluation API
will sit on (Phase 4).

\subsection{API}\label{api}

Header:
\href{../../../include/gridcalc/func/Integrate.hpp}{\passthrough{\lstinline!include/gridcalc/func/Integrate.hpp!}}.

\begin{lstlisting}[language={C++}]
namespace gridcalc::func {

struct Pairwise {};   // tag selecting pairwise recursive summation
struct Kahan {};      // tag selecting Kahan/Neumaier compensated summation

double integrate(const core::Field<double>& field, Pairwise = {}) noexcept;
double integrate(const core::Field<double>& field, Kahan)        noexcept;

}  // namespace gridcalc::func
\end{lstlisting}

Pick the reducer with a tag:

\begin{itemize}
\tightlist
\item
  \passthrough{\lstinline!integrate(f)!} or
  \passthrough{\lstinline!integrate(f, Pairwise\{\})!} ---
  \textbf{default}. Pairwise recursive summation; rounding error grows
  like \passthrough{\lstinline!O(ε log n)!}. Cheap. Use this unless you
  have a reason not to.
\item
  \passthrough{\lstinline!integrate(f, Kahan\{\})!} --- Neumaier-style
  compensated summation; \textasciitilde bit-perfect accuracy at
  \textasciitilde3-4× cost per element. Use when integrating long sums
  of values whose magnitudes vary by many orders of magnitude.
\end{itemize}

The grid's cell volume \passthrough{\lstinline!hx * hy * hz!} is read
from the input \passthrough{\lstinline!Field!}'s
\passthrough{\lstinline!Grid!} and applied by
\passthrough{\lstinline!integrate()!} itself. You do not multiply by
\passthrough{\lstinline!dV!} yourself.

Periodic wrap is irrelevant here: every grid point is sampled exactly
once, regardless of the input field's
\passthrough{\lstinline!IndexPolicy!}.

\subsection{Worked example}\label{worked-example}

Volume of a 2π cube:

\begin{lstlisting}[language={C++}]
#include <gridcalc/core/Field.hpp>
#include <gridcalc/core/Grid.hpp>
#include <gridcalc/func/Integrate.hpp>

const int N = 32;
const double L = 2.0 * M_PI;
const double h = L / N;

gridcalc::core::Grid grid(N, N, N, gridcalc::core::Vec3d(h, h, h));
gridcalc::core::Field<double> ones(grid, 1.0);

double V = gridcalc::func::integrate(ones);
// V == L*L*L == (2*pi)^3, to machine precision.
\end{lstlisting}

A periodic sinusoid integrates to zero (every axis-orthogonal trig
component sums exactly to zero on equispaced periodic samples):

\begin{lstlisting}[language={C++}]
gridcalc::core::Field<double> psi(grid, [](double x, double y, double z) {
  return std::cos(x) * std::sin(2.0 * y) * std::sin(3.0 * z);
});
double I = gridcalc::func::integrate(psi);   // I ~ 0  (within ~1e-15)
\end{lstlisting}

\subsection{\texorpdfstring{What this does \emph{not} do
(yet)}{What this does not do (yet)}}\label{what-this-does-not-do-yet}

\begin{itemize}
\tightlist
\item
  \textbf{No callable-integrand overload.}
  \passthrough{\lstinline!∫ f(ψ, ∇ψ, ∇²ψ) dV!} for user-supplied
  \passthrough{\lstinline!f!} is the headline of Phase 4 (functional
  evaluation).
\item
  \textbf{No vector-field overload.}
  \passthrough{\lstinline!integrate(Field<Vec3d>)!} arrives when the
  Phase 11+ higher-rank functionals first need it.
\item
  \textbf{No threading.} Phase 3 is single-threaded; Phase 20
  (performance pass) parallelizes the reducer while preserving
  bit-identical output.
\item
  \textbf{No non-periodic boundary handling.} The midpoint formula
  assumes the domain is periodic. Once Phase 18 introduces
  Dirichlet/Neumann fields, endpoint corrections (Euler--Maclaurin,
  Simpson) become relevant.
\end{itemize}

For the why-this-formula-is-correct background --- derivation, error
analysis, references --- see
\href{../../developer-note/notes/phase-3-domain-integration.md}{\passthrough{\lstinline!docs/developer-note/notes/phase-3-domain-integration.md!}}.
